\documentclass[12pt,a4paper]{article}
\usepackage[utf8]{inputenc}
\usepackage{listings}
\usepackage{xcolor}
\usepackage{geometry}
\geometry{a4paper, margin=1in}
\usepackage{hyperref}

% Code listing style setup
\definecolor{codegreen}{rgb}{0,0.6,0}
\definecolor{codegray}{rgb}{0.5,0.5,0.5}
\definecolor{codepurple}{rgb}{0.58,0,0.82}
\definecolor{backcolour}{rgb}{0.95,0.95,0.92}

\lstdefinestyle{mystyle}{
    backgroundcolor=\color{backcolour},   
    commentstyle=\color{codegreen},
    keywordstyle=\color{magenta},
    numberstyle=\tiny\color{codegray},
    stringstyle=\color{codepurple},
    basicstyle=\ttfamily\footnotesize,
    breakatwhitespace=false,         
    breaklines=true,                 
    captionpos=b,                    
    keepspaces=true,                 
    numbers=left,                    
    numbersep=5pt,                  
    showspaces=false,                
    showstringspaces=false,
    showtabs=false,                  
    tabsize=2
}
\lstset{style=mystyle}

\title{\textbf{Secure Drone Communication System}\\ \large Code and Short Report}
\author{Roll Number: CS23B1059}
\date{\today}

\begin{document}

\maketitle

\section{Short Report}

\subsection{System Design}
The system uses an Object-Oriented approach with \texttt{DroneClient} and \texttt{GroundStation} acting as the communicating entities. The protocol strictly follows a multi-phase workflow:
\begin{enumerate}
    \item \textbf{Authentication Phase}: The Ground Station issues a random 32-byte challenge. The Drone responds with an HMAC-SHA256 of the challenge using a pre-shared secret (Challenge-Response mechanism).
    \item \textbf{Key Exchange Phase}: The parties establish a shared secret via Elliptic Curve Diffie-Hellman (ECDH) using SECP256R1. A 32-byte MAC key is derived from this shared secret via HKDF.
    \item \textbf{Data Protection Phase}: 
    \begin{itemize}
        \item A random AES-256 session key is generated by the Drone and encrypted using the Server's RSA public key (Hybrid Encryption).
        \item Telemetry data is encrypted using AES-CBC with a random Initialization Vector (IV).
        \item A random Nonce is generated to prevent replay attacks.
        \item The encrypted payload, keys, IV, and nonce are concatenated. The drone then signs the canonical packet with ECDSA and appends an HMAC using the derived MAC key.
    \end{itemize}
\end{enumerate}

\subsection{Explanation of Implementation}
The implementation comprises several modules for clear separation of concerns. \texttt{authentication.py} handles the Challenge-Response logic. \texttt{key\_exchange.py} manages ECDH key generation and HKDF derivation. \texttt{encryption.py} provides RSA-OAEP for exchanging the AES key and AES-CBC with PKCS7 padding for encrypting the telemetry data. \texttt{signature.py} uses ECDSA for digital signatures, and \texttt{integrity.py} ensures payload integrity via HMAC-SHA256. Replay attacks are mitigated by \texttt{replay.py}, which tracks used nonces. Everything is orchestrated by \texttt{main.py}.

\subsection{Security Analysis}
\subsubsection{What attacks are prevented?}
\begin{itemize}
    \item \textbf{Replay Attacks}: A strict nonce cache rejects any payload bearing a previously processed nonce.
    \item \textbf{Man-in-the-Middle (MitM) / Forgery}: MitM is prevented because the symmetric AES key is encrypted via the Ground Station's public RSA key, and the overall payload integrity and origin are verified via both HMAC-SHA256 and an ECDSA signature bound to the drone.
    \item \textbf{Eavesdropping (Loss of Confidentiality)}: Attackers cannot read telemetry data without the AES session key, which is securely transmitted using RSA.
\end{itemize}

\subsubsection{What are the limitations?}
\begin{itemize}
    \item The Nonce cache currently grows indefinitely in memory, which over long periods could cause a memory leak. In a real system, nonces should be combined with timestamps so older nonces can be safely pruned from memory.
    \item The initial Challenge-Response relies on a symmetric pre-shared secret, meaning if the Ground Station is compromised, the drone's authentication secret is also compromised.
\end{itemize}

\newpage
\section{Python Code}

\subsection{\texttt{main.py}}
\begin{lstlisting}[language=Python]
import base64
import json
import os

from authentication import generate_challenge, compute_response, verify_response
from key_exchange import generate_ecdh_keys, derive_shared_secret
from encryption import generate_rsa_keypair, rsa_encrypt, rsa_decrypt, aes_cbc_encrypt, aes_cbc_decrypt
from signature import generate_ecdsa_keys, ecdsa_sign, ecdsa_verify
from integrity import generate_hmac, verify_hmac
from replay import generate_nonce, validate_nonce

class DroneClient:
    def __init__(self, drone_id, pre_shared_secret):
        self.drone_id = drone_id
        self.pre_shared_secret = pre_shared_secret
        self.ecdsa_priv, self.ecdsa_pub = generate_ecdsa_keys()
        self.ecdh_priv, self.ecdh_pub = generate_ecdh_keys()
        self.aes_session_key = os.urandom(32)
        
    def respond_to_challenge(self, challenge):
        print(f"[Drone] Responding to authentication challenge...")
        return compute_response(challenge, self.pre_shared_secret)
        
    def get_public_keys(self):
        return self.ecdh_pub, self.ecdsa_pub
        
    def prepare_secure_payload(self, server_rsa_pub, server_ecdh_pub):
        print("\n[Drone] Initiating secure payload preparation...")
        
        # 1. Establish derived MAC key using ECDH
        derived_key = derive_shared_secret(self.ecdh_priv, server_ecdh_pub)
        mac_key = derived_key[:32] # 32 bytes for HMAC
        
        # 2. Encrypt AES session key with Server's RSA public key
        enc_session_key = rsa_encrypt(server_rsa_pub, self.aes_session_key)
        
        # 3. Encrypt telemetry data with AES-CBC using the AES session key
        telemetry = {
            "drone_id": self.drone_id,
            "latitude": 34.0522,
            "longitude": -118.2437,
            "speed": 65
        }
        plaintext = json.dumps(telemetry).encode()
        iv, ciphertext = aes_cbc_encrypt(self.aes_session_key, plaintext)
        
        # 4. Generate Nonce for replay protection
        nonce = generate_nonce()
        
        # 5. Pack data
        packet = {
            "drone_id": self.drone_id,
            "nonce": base64.b64encode(nonce).decode(),
            "enc_session_key": base64.b64encode(enc_session_key).decode(),
            "iv": base64.b64encode(iv).decode(),
            "ciphertext": base64.b64encode(ciphertext).decode()
        }
        
        # 6. Integrity and Authentication of payload
        canonical_data = json.dumps(packet, sort_keys=True).encode()
        mac = generate_hmac(mac_key, canonical_data)
        signature = ecdsa_sign(self.ecdsa_priv, canonical_data)
        
        packet["mac"] = base64.b64encode(mac).decode()
        packet["signature"] = base64.b64encode(signature).decode()
        
        print("[Drone] Payload securely prepared and signed.")
        return packet


class GroundStation:
    def __init__(self, expected_drone_id, pre_shared_secret):
        self.expected_drone_id = expected_drone_id
        self.pre_shared_secret = pre_shared_secret
        self.rsa_priv, self.rsa_pub = generate_rsa_keypair()
        self.ecdh_priv, self.ecdh_pub = generate_ecdh_keys()
        
    def authenticate_drone(self, drone):
        print("\n[Ground Station] Sending authentication challenge to drone...")
        challenge = generate_challenge()
        response = drone.respond_to_challenge(challenge)
        
        if verify_response(challenge, response, self.pre_shared_secret):
            print("[Ground Station] Drone authentication SUCCESSFUL.")
            return True
        else:
            print("[Ground Station] Drone authentication FAILED.")
            return False
            
    def receive_payload(self, packet, drone_ecdh_pub, drone_ecdsa_pub):
        print("\n[Ground Station] Receiving secure payload...")
        
        # Extract fields
        try:
            nonce = base64.b64decode(packet.pop("nonce"))
            enc_session_key = base64.b64decode(packet.pop("enc_session_key"))
            iv = base64.b64decode(packet.pop("iv"))
            ciphertext = base64.b64decode(packet.pop("ciphertext"))
            mac = base64.b64decode(packet.pop("mac"))
            signature = base64.b64decode(packet.pop("signature"))
        except KeyError as e:
            raise ValueError(f"Missing field in packet: {e}")
            
        # Reconstruct canonical data for verification
        packet["nonce"] = base64.b64encode(nonce).decode()
        packet["enc_session_key"] = base64.b64encode(enc_session_key).decode()
        packet["iv"] = base64.b64encode(iv).decode()
        packet["ciphertext"] = base64.b64encode(ciphertext).decode()
        canonical_data = json.dumps(packet, sort_keys=True).encode()
        
        # 1. Replay Protection check
        print("[Ground Station] Checking for replay attacks...")
        validate_nonce(nonce)
        print("[Ground Station] Nonce is fresh.")
        
        # 2. Digital Signature check
        print("[Ground Station] Verifying ECDSA signature...")
        ecdsa_verify(drone_ecdsa_pub, canonical_data, signature)
        print("[Ground Station] Signature verified.")
        
        # 3. Message Integrity check
        print("[Ground Station] Verifying message integrity (HMAC)...")
        derived_key = derive_shared_secret(self.ecdh_priv, drone_ecdh_pub)
        mac_key = derived_key[:32]
        if not verify_hmac(mac_key, canonical_data, mac):
            raise ValueError("MAC verification failed!")
        print("[Ground Station] MAC verified.")
        
        # 4. Decrypt AES session key using RSA
        print("[Ground Station] Decrypting AES session key via RSA-OAEP...")
        aes_session_key = rsa_decrypt(self.rsa_priv, enc_session_key)
        
        # 5. Decrypt Payload
        print("[Ground Station] Decrypting telemetry data using AES-CBC...")
        plaintext = aes_cbc_decrypt(aes_session_key, iv, ciphertext)
        
        print("\n[Ground Station] >>> SECURE DATA RECOVERED:")
        print(json.loads(plaintext.decode()))
        return True


def simulate():
    print("=== SECURE DRONE COMMUNICATION PROTOCOL ===")
    
    # Pre-shared secret for HMAC Challenge-Response
    SHARED_SECRET = b"super_secret_drone_auth_key_99"
    DRONE_ID = "DRN-Alpha-1"
    
    drone = DroneClient(DRONE_ID, SHARED_SECRET)
    station = GroundStation(DRONE_ID, SHARED_SECRET)
    
    # Step 1: Authentication
    if not station.authenticate_drone(drone):
        print("Aborting communication.")
        return
        
    # Step 2: Key Exchange & Data Preparation
    drone_ecdh_pub, drone_ecdsa_pub = drone.get_public_keys()
    
    packet = drone.prepare_secure_payload(station.rsa_pub, station.ecdh_pub)
    
    # Step 3: Ground Station Processes Packet
    try:
        # Create a deep copy to prevent mutation issues in the simulation
        station.receive_payload(json.loads(json.dumps(packet)), drone_ecdh_pub, drone_ecdsa_pub)
    except Exception as e:
        print(f"FAILED TO PROCESS PACKET: {e}")
        
    # Step 4: Bonus - Replay Attack Simulation
    print("\n\n=== BONUS: REPLAY ATTACK SIMULATION ===")
    print("An attacker intercepts the packet and tries to resend it...")
    try:
        station.receive_payload(json.loads(json.dumps(packet)), drone_ecdh_pub, drone_ecdsa_pub)
        print("CRITICAL SECURITY FAILURE: Replay attack succeeded!")
    except ValueError as e:
        print(f"SUCCESS: Replay attack prevented! Reason: {e}")


if __name__ == "__main__":
    simulate()

\end{lstlisting}

\subsection{\texttt{authentication.py}}
\begin{lstlisting}[language=Python]
import os
import hmac
import hashlib

def generate_challenge():
    """Server generates a random challenge."""
    return os.urandom(32)

def compute_response(challenge, pre_shared_secret):
    """Client computes response using a pre-shared secret key."""
    return hmac.new(pre_shared_secret, challenge, hashlib.sha256).digest()

def verify_response(challenge, response, pre_shared_secret):
    """Server verifies the client's response."""
    expected = compute_response(challenge, pre_shared_secret)
    return hmac.compare_digest(expected, response)
\end{lstlisting}

\subsection{\texttt{key\_exchange.py}}
\begin{lstlisting}[language=Python]
from cryptography.hazmat.primitives.asymmetric import ec
from cryptography.hazmat.primitives import hashes
from cryptography.hazmat.primitives.kdf.hkdf import HKDF

def generate_ecdh_keys():
    """Generates an Elliptic Curve key pair for ECDH."""
    private_key = ec.generate_private_key(ec.SECP256R1())
    public_key = private_key.public_key()
    return private_key, public_key

def derive_shared_secret(private_key, peer_public_key):
    """Derives a shared secret using ECDH and HKDF."""
    shared_key = private_key.exchange(ec.ECDH(), peer_public_key)
    
    # Use HKDF to expand the shared key into a 32-byte key
    derived_key = HKDF(
        algorithm=hashes.SHA256(),
        length=32,
        salt=None,
        info=b'drone-secure-ecdh-handshake',
    ).derive(shared_key)
    
    return derived_key

\end{lstlisting}

\subsection{\texttt{encryption.py}}
\begin{lstlisting}[language=Python]
import os
from cryptography.hazmat.primitives.ciphers import Cipher, algorithms, modes
from cryptography.hazmat.primitives import padding as sym_padding
from cryptography.hazmat.primitives.asymmetric import rsa, padding as asym_padding
from cryptography.hazmat.primitives import hashes

def generate_rsa_keypair():
    """Generate RSA keys for secure AES key exchange."""
    private_key = rsa.generate_private_key(public_exponent=65537, key_size=2048)
    return private_key, private_key.public_key()

def rsa_encrypt(public_key, plaintext):
    """Encrypt using RSA-OAEP."""
    return public_key.encrypt(
        plaintext,
        asym_padding.OAEP(
            mgf=asym_padding.MGF1(algorithm=hashes.SHA256()),
            algorithm=hashes.SHA256(),
            label=None
        )
    )

def rsa_decrypt(private_key, ciphertext):
    """Decrypt using RSA-OAEP."""
    return private_key.decrypt(
        ciphertext,
        asym_padding.OAEP(
            mgf=asym_padding.MGF1(algorithm=hashes.SHA256()),
            algorithm=hashes.SHA256(),
            label=None
        )
    )

def aes_cbc_encrypt(key, plaintext_bytes):
    """Encrypt data using AES-CBC with PKCS7 padding."""
    iv = os.urandom(16)
    cipher = Cipher(algorithms.AES(key), modes.CBC(iv))
    encryptor = cipher.encryptor()
    
    padder = sym_padding.PKCS7(algorithms.AES.block_size).padder()
    padded_data = padder.update(plaintext_bytes) + padder.finalize()
    
    ciphertext = encryptor.update(padded_data) + encryptor.finalize()
    return iv, ciphertext

def aes_cbc_decrypt(key, iv, ciphertext):
    """Decrypt AES-CBC data and remove PKCS7 padding."""
    cipher = Cipher(algorithms.AES(key), modes.CBC(iv))
    decryptor = cipher.decryptor()
    padded_data = decryptor.update(ciphertext) + decryptor.finalize()
    
    unpadder = sym_padding.PKCS7(algorithms.AES.block_size).unpadder()
    plaintext_bytes = unpadder.update(padded_data) + unpadder.finalize()
    return plaintext_bytes

\end{lstlisting}

\subsection{\texttt{signature.py}}
\begin{lstlisting}[language=Python]
from cryptography.hazmat.primitives.asymmetric import ec
from cryptography.hazmat.primitives import hashes

def generate_ecdsa_keys():
    """Generates an Elliptic Curve key pair for ECDSA signatures."""
    private_key = ec.generate_private_key(ec.SECP256R1())
    public_key = private_key.public_key()
    return private_key, public_key

def ecdsa_sign(private_key, message):
    """Sign a message using ECDSA with SHA-256."""
    return private_key.sign(message, ec.ECDSA(hashes.SHA256()))

def ecdsa_verify(public_key, message, signature):
    """Verify an ECDSA signature. Raises an exception if invalid."""
    public_key.verify(signature, message, ec.ECDSA(hashes.SHA256()))
\end{lstlisting}

\subsection{\texttt{integrity.py}}
\begin{lstlisting}[language=Python]
import hmac
import hashlib

def generate_hmac(key, message):
    """Generate HMAC-SHA256."""
    return hmac.new(key, message, hashlib.sha256).digest()

def verify_hmac(key, message, mac_to_verify):
    """Verify HMAC-SHA256."""
    expected_mac = generate_hmac(key, message)
    return hmac.compare_digest(expected_mac, mac_to_verify)
\end{lstlisting}

\subsection{\texttt{replay.py}}
\begin{lstlisting}[language=Python]
import os

class NonceValidator:
    def __init__(self):
        self._nonce_cache = set()

    def generate_nonce(self):
        """Generate a random 16-byte nonce."""
        return os.urandom(16)

    def validate_nonce(self, nonce):
        """Check if nonce has been used before to prevent replay attacks."""
        if nonce in self._nonce_cache:
            raise ValueError("Replay Attack Detected: Nonce has already been used.")
        self._nonce_cache.add(nonce)

    def clear(self):
        """Clear cache (for testing purposes)."""
        self._nonce_cache.clear()

# Global default validator
default_validator = NonceValidator()

def generate_nonce():
    return default_validator.generate_nonce()

def validate_nonce(nonce):
    default_validator.validate_nonce(nonce)

\end{lstlisting}

\end{document}
